%% kapa: exam=1 round=30 year=2019 period=2 form=A subject=appraisal_law subject_ko=감정평가관계법규 seq=40 number=40
%% source: Q-Net 감정평가사 기출문제 (cst00309) artlSeq=5205205 file=extracted/1008/5205205/2019년 제30회 감정평가사 1차시험 문제지/2019년 제30회 감정평가사 1차시험 문제지/2019년도 제30회 감정평가사 1차 2교시 A형.hwp
%% conversion: kordoc 4.24.51 -> tools/gen1.py (auto); review=proofread-sonnet; flags=none
\begin{question}{40}
도시 및 주거환경정비법령상 정비사업의 시행에 관한 설명으로 옳지 않은 것은?
\begin{choices}
  \item 재건축사업은 조합이 조합원의 과반수의 동의를 받아 시장ㆍ군수등과 공동으로 시행할 수 있다.
  \item 토지등소유자가 20인 미만인 경우에는 토지등소유자가 직접 재개발사업을 시행할 수 없다.
  \item 조합설립추진위원회도 개략적인 정비사업 시행계획서를 작성할 수 있다.
  \item 재개발사업은 정비구역에서 인가받은 관리처분계획에 따라 건축물을 건설하여 공급하거나 환지로 공급하는 방법으로 한다.
  \item 조합이 사업시행자인 경우 시장ㆍ군수등은 특별한 사유가 없으면 사업시행계획서의 제출이 있은 날부터 60일 이내에 인가 여부를 결정하여 사업시행자에게 통보하여야 한다.
\end{choices}
\end{question}
\answer{2}
